% sec-03: Set A, applications 01 to 05.  Figures 5 (plantuml use case) and
% 8 (mermaid decision flowchart).
\section{Set A: The Surgical Perspective}

\subsection{What Leads in Set A}

In the five surgical applications the operation leads. Daraxonrasib is present
in all five, but it enters as the perioperative arm the operation is timed
around rather than as the headline, and the endpoint table is led by conversion,
anastomotic integrity, operative time, estimated blood loss, R0 margin status,
and 90-day morbidity graded to the international definition
\cite{bassi2017isgps}.

That ordering is not cosmetic. A reviewer at a mechanism office asks whether the
hard part is buildable; a reviewer at a clinical programme asks whether the hard
part is survivable. Set A answers the first, and the hard part it names is the
operation with a language model in the room.

\subsection{Who Is Authorised to Do What}

The single most common objection to a language model in an operating theatre is
that the boundary is asserted rather than specified. Applications 02 and 10 each
carry a three-actor slice of the authority question. Figure 5 gives the whole
set, including the two actors that appear in no application figure: the research
pharmacist and the institutional review board.

\begin{appfloat}[!tb]
\begin{appfig}[plantuml-type / use case]
% Six human actors in the left column at x = 0, pitch 1.55.  Eleven use cases in
% two columns inside the boundary at x = 5.0 and 8.6, pitch 1.30.  The advisory
% model sits alone on the right at x = 12.6 so its associations are read against
% the human ones.  Four struck links mark the paths it does not have.
\node[umlpkg,minimum width=7.6cm,minimum height=8.4cm] (sys) at (6.8,-3.4) {};
\node[umlpkgtab,anchor=north west] at ([xshift=1mm,yshift=-1mm]sys.north west) {Trial system boundary};
\node[umlusekey,text width=22mm] (u1) at (5.0,-0.4) {Approve a robotic motion};
\node[umlusekey,text width=22mm] (u2) at (5.0,-1.7) {Order conversion to open};
\node[umlusekey,text width=22mm] (u3) at (5.0,-3.0) {Trigger arm or system stop};
\node[umlusecase,text width=22mm] (u4) at (5.0,-4.3) {Exchange or dock an instrument};
\node[umlusecase,text width=22mm] (u5) at (5.0,-5.6) {Set or hold the drug dose};
\node[umlusecase,text width=22mm] (u6) at (8.6,-0.4) {Adjudicate attribution};
\node[umlusecase,text width=22mm] (u7) at (8.6,-1.7) {Halt the cohort};
\node[umlusecase,text width=22mm] (u8) at (8.6,-3.0) {Suspend the study};
\node[umlusegray,text width=22mm] (u9) at (8.6,-4.3) {Render advisory text};
\node[umlusegray,text width=22mm] (u10) at (8.6,-5.6) {Write to the audit log};
\node[umlusegray,text width=22mm] (u11) at (8.6,-6.9) {Read telemetry, read only};
\umlactor{a1}{0}{-0.3}{Operating surgeon}
\umlactor{a2}{0}{-2.0}{Bedside assistant}
\umlactor{a3}{0}{-3.7}{GI medical oncologist}
\umlactor{a4}{0}{-5.4}{Research pharmacist}
\umlactor{a5}{0}{-7.1}{Independent safety monitor}
\umlactor{a6}{2.3}{-7.7}{IRB}
\umlactor{a7}{12.6}{-4.6}{On-premises advisory model}
\draw[umlassoc] (a1) -- (u1.west);
\draw[umlassoc] (a1) -- (u2.west);
\draw[umlassoc] (a1) -- (u3.west);
\draw[umlassoc] (a2) -- (u4.west);
\draw[umlassoc] (a3) -- (u5.west);
\draw[umlassoc] (a3) -- (u6.west);
\draw[umlassoc] (a4) -- (u5.west);
\draw[umlassoc] (a5) -- (u7.west);
\draw[umlassoc] (a6) -- (u8.west);
\draw[umlassoc] (a7) -- (u9.east);
\draw[umlassoc] (a7) -- (u10.east);
\draw[umlassoc] (a7) -- (u11.east);
\foreach \tx/\ty in {11.1/-0.9, 11.1/-2.0, 11.1/-3.1, 11.1/-6.2}{\pxmark{\tx}{\ty}}
\node[umlnote,anchor=north west,text width=34mm] (n1) at (10.6,0.4)
  {No association exists from the model to motion approval, conversion, stop, or
   dose. The four denials are structural, not procedural.};
\end{appfig}
\vspace{-0.7cm}
\figcaption{Seven actors and the eleven actions each may take. Every action has
at least one human association; the advisory model has exactly one, and the
struck links are the four it is denied by construction.}
\end{appfloat}

Three fills separate three authority classes. Corporate Blue marks the three
actions only the operating surgeon takes. Pale blue marks the shared clinical
and oversight actions. Gray marks the three actions the model may take, all of
which produce text or read data and none of which moves anything. A reader who
ignores the labels entirely still reads the structure correctly.

\subsection{Applications 01 to 05}

\textbf{01, NIH Director's Pioneer Award.} The unit of risk is the individual
rather than the aim. The bet, that a targeted RAS(ON) inhibitor and a robotic
operation can be evaluated as one governed system, does not survive
decomposition into three independent project aims, because the drug schedule,
the operative plan, and the advisory boundary are only interesting together. The
ask is \$700,000 per year for five years, with a stated go or no-go condition at
each of three gates.

\textbf{02, ARPA-H.} Written as a gated programme rather than a project
narrative, because the report endorses the ARPA model on the specific ground
that a program manager can end a bet. Three gates at months 12, 21, and 33, each
carrying a numeric kill criterion checkable from programme data: IND not active
or measured stop latency above specification; any unexpected grade 4 event
attributed to the device or the advisory path; fewer than two evaluable
participants at any dose level.

\textbf{03, NSF TIP X-Labs.} Argues organization type rather than money. The
programme needs five activities held together, operating under an active IND,
maintaining a verification harness, publishing negative operative results,
holding the regulatory package as a maintained product, and releasing code and
data openly. No existing performer type covers all five, and that gap is what
the report says X-Labs was created to fill.

\textbf{04, DOE Genesis Mission.} The FY 2028 annex defines the Robotics
national mission as ``general-purpose autonomous systems capable of dexterous
manipulation, mobility, and reliable operation in real-world environments, to
initiate the era of physical AI-driven scientific discovery''
\cite{kratsiosvought2026fy2028priorities}. A human operating room is the hardest
available instance of that sentence. Three artifacts are offered to the mission
rather than to the clinic: a trust-boundary topology, a closed loop with one
measured human gate, and open data formats \cite{eo14363genesismission}.

\textbf{05, NIH SEED SBIR.} The only application in which the commercial product
carries the argument. The deliverable is not the robot and not the drug but the
governance package that makes an on-premises model usable in a regulated
theatre: a logging schema, a bench stop-latency method, and a regulatory dossier
a qualified programme can adopt without rebuilding either.

\subsection{What a Reviewer Actually Asks}

\begin{appfloat}[!tb]
\begin{appfig}[mermaid-type / flowchart with decisions]
% Five decision diamonds on one rank at y = 0, pitch 2.9, aspect 2.1.  The
% section that answers each sits directly above it at y = 1.9, so the pairing is
% read from position and the dashed edges only confirm it.
\foreach \i/\x/\q in {1/0/{1. Is the person credible?}, 2/2.9/{2. Is the science plausible?},
  3/5.8/{3. Is the site reachable?}, 4/8.7/{4. Is the budget proportionate?},
  5/11.6/{5. Is the claim bounded?}}{
  \node[mmdec,text width=21mm] (g\i) at (\x,0) {\q};}
\foreach \i/\x/\s in {1/0/{\S1: 14 deposited works}, 2/2.9/{\S3: QSP, twin, RASolute 302},
  3/5.8/{\S5: Moores feasibility route}, 4/8.7/{\S5: \$700K per year, gated},
  5/11.6/{back matter: scope of claims}}{
  \node[mmsoft,text width=25mm] (s\i) at (\x,1.9) {\s};
  \draw[mmedged] (s\i) -- (g\i);}
\foreach \i/\j in {1/2, 2/3, 3/4, 4/5}{
  \draw[mmedgeb] (g\i) -- node[mmlabel,pos=0.5] {yes} (g\j);}
\node[mmgray,text width=24mm] (e1) at (1.45,-2.2) {Review ends};
\node[mmgray,text width=24mm] (e2) at (8.7,-2.2) {Ask is revised};
\node[mmgoal,text width=24mm] (fund) at (11.6,-2.2) {Proceed};
\draw[mmedgek] (g1.south) -- node[mmlabel,pos=0.55] {no} (e1.north west);
\draw[mmedgek] (g2.south) -- node[mmlabel,pos=0.55] {no} (e1.north east);
\draw[mmedge] (g3.south) -- node[mmlabel,pos=0.55] {no} (e2.north west);
\draw[mmedge] (g4.south) -- node[mmlabel,pos=0.55] {no} (e2.north);
\draw[mmedgeb] (g5.south) -- node[mmlabel,pos=0.5] {yes} (fund.north);
\end{appfig}
\vspace{-0.7cm}
\figcaption{The five questions a reviewer asks in order, and the section of the
application that answers each. A no at gate 1 or gate 2 ends the review; a no at
gates 3 to 5 changes the ask rather than ending it.}
\end{appfloat}

The asymmetry between the first two gates and the last three is why the
applications are ordered as they are. Credibility and plausibility are settled
in the first two sections of every attachment, before any ask appears, because a
reviewer who is unconvinced there never reaches the budget.

\begin{apptable}
\begin{tabularx}{\textwidth}{>{\raggedright\arraybackslash}p{0.7cm} >{\raggedright\arraybackslash}p{1.5cm} >{\raggedright\arraybackslash}p{4.6cm} >{\raggedright\arraybackslash}X}
\headrow \thead{No} & \thead{Figures} & \thead{Types used} & \thead{Type deliberately not used, and why}\\
\midrule
01 & 3 & mermaid, d2, graphviz & No plantuml or diagrams: no authority question, no topology question \\
02 & 4 & mermaid, plantuml, d2, graphviz & No diagrams: nothing is deployed across a boundary in this argument \\
03 & 3 & d2, diagrams, mermaid & No plantuml or graphviz: the argument is organization type, not authority or dependency \\
04 & 4 & diagrams, mermaid, graphviz, d2 & No plantuml: the guards are stated in prose and drawn in the summary paper instead \\
05 & 3 & mermaid, graphviz, d2 & No plantuml or diagrams: a phased commercial argument needs neither \\
\end{tabularx}
\end{apptable}
\tabcap{Figure inventory for Set A, including the platform each application
deliberately did not use. Recording the omission is what keeps the type choice
a decision rather than a default.}
